\documentclass[10pt,a4paper]{amsart}
\usepackage[margin=19mm]{geometry}
\usepackage{amsmath,amssymb,mathtools}
\usepackage[scr=rsfs,cal=euler]{mathalfa}
\usepackage{xfrac}
\usepackage[unicode,hidelinks,bookmarks=false]{hyperref}
\newcommand{\ie}{i.e.\ }
\newcommand{\cf}{cf.\ }
\newcommand{\resp}{respectively\ }
\newcommand{\Subsection}{\subsection}
\providecommand{\nobreakdash}{\nobreak\hbox{-}}
\setlength{\emergencystretch}{2em}
\multlinegap0pt
\usepackage{bm}
\usepackage{bbm}
\usepackage{dsfont}
\usepackage{xspace}
\usepackage{cancel}
\usepackage{mathtools}
\usepackage{cleveref}
\usepackage{amsthm}
\makeatletter
\@ifundefined{theorem}{\newtheorem{theorem}{Theorem}}{}
\@ifundefined{lemma}{\newtheorem{lemma}[theorem]{Lemma}}{}
\@ifundefined{proposition}{\newtheorem{proposition}[theorem]{Proposition}}{}
\makeatother
\newcommand{\N}{\mathbb{N}}
\title[On subshifts with low maximal pattern complexity]{On subshifts with low maximal pattern complexity}
\author{Anh N. Le, Ronnie Pavlov,, Casey Schlortt}
\date{Source: arXiv:2508.13420v1 (CC BY 4.0)}
\begin{document}
\maketitle

\noindent\emph{Source and licence.} On subshifts with low maximal pattern complexity. Version arXiv:2508.13420v1 (CC BY 4.0), distributed under the Creative Commons Attribution 4.0 International licence, as recorded in the arXiv licence metadata for this exact version and shown on the abstract page https://arxiv.org/abs/2508.13420v1. Full rights notice: licenses/arxiv-2508.13420-CC-BY-4.0.txt.\par\smallskip

\noindent\textit{Editorial scope.} Counted unit: the Proposition whose source label is \texttt{prop:w-constant} in the article (source0.tex lines 1405--1407, source label \texttt{prop:w-constant}), delivered together with the article's own complete original proof of it (source0.tex lines 1408--1471, one \texttt{proof} environment of 64 lines). No mathematics is written, repaired or re-derived: statement and proof are the article's own text line for line. Required context retained uncounted: lem:long\_0\_and\_1\_not\_pattern\_Sturmian (lines 1345--1347). Every definition, display and label of those lines is retained unchanged. Omitted from this package and not redistributed: 1--1344, 1348--1404, 1472--1564. Nothing inside a retained excerpt is omitted or altered: every retained body is delivered exactly as the article's lines are written, so no omission receipt of any kind accompanies this package. Citations: the retained excerpts carry no citation command at all, so no bibliography entry of the article is transcribed and no reference is redistributed. Reference adaptation: none was needed and none was made; every reference inside the delivered unit resolves inside it, so no unresolved reference or citation is produced. Printed slips kept literal: no printed word, symbol, hypothesis or number of the counted statement or of its proof was altered. Layout and class: the article's own class is \texttt{amsart} with packages including amssymb, amscd, inputenc, amsmath, amsfonts, faktor, latexsym, verbatim, graphicx, epsfig, enumitem, hyperref, cleveref, geometry; the wrapper is the lane's pinned \texttt{amsart} editorial wrapper at 10pt on A4, which restates the theorem-like environments the retained text needs and adds the article's own macro definitions that the retained text needs (\texttt{\textbackslash N}), with extra wrapper packages (bm, bbm, dsfont, xspace, cancel, mathtools, cleveref). The delivered numbering is the unit's own. Assets: no retained span contains a figure environment, a graphics-inclusion command, a PDF-page inclusion, a table or a TikZ picture; no image file is copied into this package, the package declares no asset dependency and no image file is redistributed with it. Licence: version arXiv:2508.13420v1 of this article is distributed under the Creative Commons Attribution 4.0 International licence, as recorded in the arXiv licence metadata for this exact version and shown on the abstract page https://arxiv.org/abs/2508.13420v1. A full rights notice is delivered at licenses/arxiv-2508.13420-CC-BY-4.0.txt. Characters: every retained excerpt is pure ASCII, so the delivered engine, XeLaTeX with its default Latin Modern text font, reports no missing character.
\begin{lemma}\label{lem:long_0_and_1_not_pattern_Sturmian}
If $x \in \{0, 1\}^{\N_0}$ has arbitrarily long blocks of consecutive $0$s and arbitrarily long blocks of consecutive $1$s, then $x$ is not pattern Sturmian. 
\end{lemma}

\begin{proposition}\label{prop:w-constant}
Let $x \in \{0, 1\}^{\N_0}$ be such that $x$ differs from the infinite sequence $w w w \ldots$ on a set of Banach density zero for some finite word $w$. If $w$ is not a constant word, then $x$ is not pattern Sturmian. 
\end{proposition}

\begin{proof}
Assume that $x$ and $w$ are as in the statement, and assume that $w$ contains both $0$ and $1$. 
Let $k = |w|$, the length of $w$. Consider the sequences $(x^{(j)}(n))_{n \in \N_0} = (x(kn + j))_{n \in \N_0}$ for $j \in \{0, \ldots, k - 1\}$. For each $j$, $x^{(j)}$ contains arbitrarily long blocks of consecutive $0$s or $1$s or both. However, the last possibility is ruled out by \cref{lem:long_0_and_1_not_pattern_Sturmian}, since $p^*_x(n) \geq p^*_{x^{(j)}}(n)$. 
Since $w$ contains both $0$ and $1$, there are $j_0, j_1 \in \{0, \ldots, k-1\}$ that $x^{(j_0)}$ contains arbitrarily long blocks of $0$s and $x^{(j_1)}$ contains arbitrarily long blocks of $1$s. Because $x$ is nonperiodic, we can choose $j_0, j_1$ so that at least one of $x^{(j_0)}, x^{(j_1)}$ is nonperiodic. Without loss of generality (and by switching $0$ and $1$ if necessary, which does not change any hypothesis or conclusion of the proposition), assume $x^{(j_0)}$ is nonperiodic. 

We will show that there exists a $3$-window $\tau$ such that 
\[
    |L_{x^{(j_0)}}(\tau) \cup L_{x^{(j_1)}}(\tau)| \geq 7.
\]
It then easily follows that $p_{x}^*(3) \geq 7$, implying that $x$ is not pattern Sturmian.

For any $3$-window $\tau$, $L_{\beta_0}(\tau)$ always contains $000$ and $L_{\beta_1}(\tau)$ always contains $111$. Moreover, if $n$ is the location of an $1$ before entering a long block of $0$s, then $x(n + \tau) = 100$. Thus $L_{\beta_0}(\tau)$ always contains $100$, and a similar argument shows it always contains $001$.


Let $S = \{s_1 < s_2 < \ldots\}$ be the set of locations where $1$ appears in $x^{(j_0)}$ and define $g_n = s_{n+1} - s_n$. The set $S$ is not syndetic because $d^*(S) = 0$ and so there exists $n$ such that 
\begin{equation}\label{eq:condition_s_n}
    g_n < g_{n+1}.
\end{equation}
Choose the window
\begin{equation}\label{eq:choose_tau}
    \tau = \{0, g_n, g_{n+1}\}.
\end{equation}
Then 
\[
    x^{(j_0)}(s_n + \tau) = x^{(j_0)}(\{s_n, s_{n+1}, s_n + g_{n+1}\}) = 110
\]
because $s_{n+1} < s_n + g_{n+1} < s_{n+2}$. Furthermore,
\[
    x^{(j_0)}(s_{n+1} + \tau) = x^{(j_0)}(\{s_{n+1}, s_{n+1} + g_n, s_{n+2}]\}) = 101
\]
because $s_{n+1} < s_{n+1} + g_n < s_{n+2}$.

So far with the window $\tau$ in \eqref{eq:choose_tau}, we have shown $L_{x^{(j_0)}}(\tau) \cup L_{x^{(j_1)}}(\tau)$ contains $6$ words $000, 001, 100, 111, 110, 101$. It remains to find $n$ so that we pick up the extra word $010$ when sliding the window $\tau$ along $x^{(j_0)}$. For contradiction, assume there exists no such $n$. 

Fix an $n_0$ such that $g_{n_0+1} > g_{n_0}$. We claim that for all $\ell$, there exists $k \geq 2$ such that $g_{k-1} \leq g_{n_0}$ and $g_{k} > \ell$. For sufficiently large $k$ (more specifically, if $s_k > g_{n_0}$), we have
\begin{equation}\label{eq:layout}
    x^{(j_0)}(s_k - g_n + \tau) = x^{(j_0)}(\{s_k - g_n, s_k, s_k + (g_{n+1} - g_{n})\}) = u1v
\end{equation}
where $u, v \in \{0, 1\}$. By our contradiction assumption, $u = 1$ or $v = 1$. If $s_k$ is the location of a $1$ right before a very large block of $0$s (say the length of this block is larger than $\ell$), then $x^{(j_0)}(s_k + (g_{n_0+1} - g_{n_0})) = 0$ and so this forces $x^{(j_0)}(s_k - g_{n_0}) = 1$. Since $s_{k-1}$ is the location of the last $1$ before $s_k$, we have 
\[
    g_{k-1} = s_k - s_{k-1} \leq s_k - (s_k - g_{n_0}) = g_{n_0}.
\]
On the other hand, since the digit $1$ at the location $s_k$ is followed by an $\ell$-block of $0$s, $g_k = s_{k+1} - s_k > \ell$. Our claim follows.


Let $\ell$ be arbitrary and fix $k$ so that $g_{k-1} \leq g_{n_0}$ and $g_k > \ell g_{n_0} \geq \ell g_{k-1}$.
We will show that $x^{(j_0)}$ contains a sequence of $\ell + 1$ $1$ symbols where the gaps between consecutive $1$s are bounded by $g_{n_0}$. Since $\ell$ is arbitrary this will imply that the upper Banach density of $1$ in $x^{(j_0)}$ is at least $1/g_{n_0}$ and this contradicts the assumption that this density is zero. Thus we will be done. 


Let $s_t$ be the location of an $1$ before a block of consecutive $0$s of length greater than $g_k - g_{k-1}$.
Looking at \eqref{eq:layout}, because 
\[
    x^{(j_0)}(s_t) = 1 \text{ and } x^{(j_0)}(s_t + g_{k} - g_{k-1}) = 0,
\]
it must be that $x^{(j_0)}(s_t - g_{k-1}) = 1$. Now let $s_t - g_{k-1}$ plays the role of $s_t$. We then have 
\[
    x^{(j_0)}(s_t - g_{k-1}) = 1 \text{ and } x^{(j_0)}(s_t - g_{k-1} + (g_k - g_{k-1}))= x^{(j_0)}(s_t + g_k - 2 g_{k-1}) = 0
\]
force $x^{(j_0)}(s_t - 2 g_{k-1}) = 1$. Continuing in this way, we obtain $1$ at the following $\ell + 1$ locations in $x^{(j_0)}$
\[
    s_t, s_t - g_{k-1}, \ldots, s_t - \ell g_{k-1}.
\]
At each step, we use the fact that $g_k > \ell g_{k-1}$ to make sure $s_t + g_k - \ell g_{k-1} > s_t$ and so $x^{(j_0)}(s_t + g_k - \ell g_{k_1}) = 0$. We have arrived at a contradiction, and so $x$ is not pattern Sturmian, completing the proof.
\end{proof}

\end{document}
